\documentclass[../../main.tex]{subfiles}
\begin{document}

\begin{center}
  {\fontsize{48pt}{58pt}\selectfont \textbf{1994}\normalsize\textbf{年度}}
\end{center}

\begin{table}[htb]
  \centering
  \caption{出題テーマと難易度}
  \begin{tabular}{|c|p{9cm}|c|}
    \hline
    \textbf{問題番号} & \textbf{テーマ} & \textbf{難易度} \\
    \hline
    第1問 & 絶対値を含む不等式と等号成立 & \\
    \hline
    第2問 & 行列のべきと合同式 & \\
    \hline
    第3問 & 円の接線と正三角形の頂点の軌跡 & \\
    \hline
    第4問 & 3人の対戦と優勝確率 & \\
    \hline
    第5問 & 円の弦が通過する部分の最大面積 & \\
    \hline
    第6問 & 対数の積分の漸化式と不等式 & \\
    \hline
  \end{tabular}
\end{table}

\begin{table}[htb]
  \centering
  \caption{解答}
  \renewcommand{\arraystretch}{1.4}
  \begin{tabular}{|c|c|p{9cm}|}
    \hline
    \textbf{問題番号} & \textbf{小問} & \textbf{答え} \\
    \hline
    第1問 & -- & 不等式は成立し，等号は$abc=0$のとき \\
    \hline
    第2問 & (1)，(2) &  \\
    \hline
    第3問 & -- & $\sqrt{3}(a^2+b^2)-4ab-2+2(\sqrt{3}-1)(a+b)=0$，$\sqrt{3}b>a$，$\sqrt{3}a>b$ \\
    \hline
    第4問 & (1) & $1-p$ \\
    \cline{2-3}
    & (2) & $p(1-p)$ \\
    \cline{2-3}
    & (3) & $p=\dfrac{\sqrt{3}}{3}$ \\
    \hline
    第5問 & -- & 最大値$\dfrac{1}{2\pi}$ \\
    \hline
    第6問 & (1) & $I_{n+1}=e-(n+1)I_n$ \\
    \cline{2-3}
    & (2) & 証明問題 \\
    \hline
  \end{tabular}
\end{table}

\paragraph{試験形態}
6問．試験時間は要確認．

\end{document}
